\section{Compiling saturated bundles without expanding their genus}
\label{sec:relative-geometry}

A parallelity component with every base boundary circle attached to the core
differs fundamentally from one with an unattached boundary interval.  The
incoming hierarchy-kernel report proves that a proper-boundary attachment
adds only relative one-handles.  It also shows that a fully attached bordered
base requires at most one additional two-handle, but leaves that handle's
attaching curve uncompiled~\cite{incoming-hierarchy-kernels}.  The result of
this section supplies the geometric intersection data needed to use that
remaining handle in the maintained cancellation model.

The new information is more precise than a relation in a fundamental-group
presentation.  We identify the unsigned intersections of an \emph{embedded}
residual attaching curve with the original core cocores.  Its potentially
numerous additional one-handles can then be omitted from the cancellation
matrix, while their number is retained exactly.  The resulting certificate
determines an unmarked handlebody census.  It does not specify the position
of every inherited boundary-pattern curve after the cancellations.

\subsection{The geometric input and its source product charts}
\label{geom:input}

Let $K$ be a compact orientable three-manifold supplied with an actual
handle decomposition containing zero-, one-, and two-handles and no
three-handles.  Write $H$ for the union of its zero- and one-handles.
Thus $H$ is a disjoint union of handlebodies; we do not initially assert
that the components of $K$ are handlebodies.  For an original one-handle
$h$, let $D_h$ denote its cocore disc.  The original two-handle attaching
curves are embedded, pairwise disjoint curves on $\partial H$.

Suppose that orientable total interval bundles
\[
 p_i:E_i\longrightarrow B_i
\]
are attached to $K$, where every $B_i$ is a connected compact surface
with nonempty boundary.  Their complete intersections with the core are
$p_i^{-1}(A_i)$, for compact one-dimensional submanifolds
$A_i\subseteq\partial B_i$.  Distinct bundle interiors are disjoint,
and the interfaces have compatible corner collars and no unrecorded
intersections.  Adjacent source ports have already been amalgamated at
their actual common corners, so the connected components of $A_i$ are
genuine maximal closed intervals or whole boundary circles.

For an attached component, $A_i\ne\varnothing$.  Put
\[
 a_i=\#\{\text{interval components of }A_i\},\qquad
 c_i=\#\{\text{circle components of }A_i\},\qquad
 b_i=|\pi_0(\partial B_i)|.
\]
We call the component \emph{saturated} if $A_i=\partial B_i$ and set
$\varepsilon_i=1$ in that case; otherwise set $\varepsilon_i=0$.
For a genuine boundary submanifold this is equivalent to
$c_i=b_i$, in which case necessarily $a_i=0$.

\begin{definition}[Compatibility with the source cocores]
\label{geom:def-compatible}
The attaching charts are compatible with the original core cocores if,
after the fixed compatible rounding and amalgamation of their specified
corner contacts, their entire port rectangles and collars lie in the
surviving part of $\partial H$ outside the original two-handle attaching
regions, and every core-facing port has one of the following product
descriptions:
\begin{enumerate}
\item A core-zero port lies in the boundary of its specified zero-handle
and misses all $\partial D_h$.
\item A core-one port is a rectangle whose base-arc coordinate runs
between the two ends of its specified one-handle and whose other
coordinate is the interval-bundle fibre.  It meets $\partial D_h$, for
that one-handle $h$, in one complete fibre, and misses every other
$\partial D_{h'}$.
\end{enumerate}
The statements concern the entire rectangles, including their collars,
and not only the signed words of their midpoint paths.
\end{definition}

These are the original source charts used by the maintained disk-bundle
absorption operator~\cite{marked-parallelity-85}.  A core-one port comes
from the face-sector side of an edge-cap polygon.  The transverse
coordinate of the source one-handle runs from one tetrahedron chamber
to the other; the bundle fibre runs through the gap between parallel
normal sheets.  A middle transverse cocore therefore cuts the port in
one complete fibre.  The source tetrahedron--face flags distinguish its
two ends even for a self-pairing.  This product fact justifies
Definition~\ref{geom:def-compatible} for the native adapter.

The midpoint section is preserved by a fibre reversal.  Consequently a
whole attaching circle has a well-defined embedded midpoint curve
$\gamma_{ij}$, with actual unsigned incidence
\begin{equation}
 n_{ij,h}=|\gamma_{ij}\cap\partial D_h|
 =\#\{\text{core-one ports on that circle owned by }h\}.
 \label{geom:eq-annulus-incidence}
\end{equation}
The finite source walk also verifies that consecutive port segments
join and that the circle closes.  Neither an exponent sum nor a freely
reduced word can replace the count in
\eqref{geom:eq-annulus-incidence}.

The cocore compatibility is a substantive hypothesis.  For an arbitrary
annulus parametrization and an arbitrary meridian system, a cocore
boundary can have an additional component running parallel to the
annulus core.  Midpoint intersection numbers alone would then fail to
certify that an entire band foot avoids that boundary.  We compile the
new rows against the original source cocores, before any disc slides,
so that the verified rectangle description remains available.

\subsection{Relative bands and the contribution of a saturated base}
\label{geom:relative-counts}

We recall the relative count from the incoming report, including the
geometric reason for it.  This count is prior work within the project;
the residual incidence theorem below is the extension developed here.

\begin{lemma}[Relative band count]
\label{geom:lem-relative-count}
Let $B$ be a connected compact bordered surface and let
$\varnothing\ne A\subseteq\partial B$ have $a$ interval components
and $c$ circle components.  Set
$\varepsilon=\mathbf{1}_{A=\partial B}$.  An orientable total interval
bundle over $B$, attached along its vertical bundle over $A$, admits a
relative decomposition with
\begin{equation}
 h_1=a-\chi(B)+\varepsilon,
 \qquad h_2=\varepsilon,
 \label{geom:eq-relative-count}
\end{equation}
and no additional zero- or three-handles.  The one-handles connect the
$a+c$ initial interface collars.  In particular
\[
 \lambda=h_1-(a+c)+1=1-\chi(B)-c+\varepsilon\geq0.
\]
\end{lemma}

\begin{proof}
First suppose that $A\ne\partial B$.  Triangulate the pair $(B,A)$ and
choose a boundary edge outside $A$.  Root a spanning tree of the dual
triangle graph at a triangle incident to that edge.  Collapse this
triangle across the chosen boundary edge, and subsequently collapse
each triangle across its edge to its already removed parent.  These
are free edges when used.  No chosen edge lies in $A$, so the result
is a connected graph containing $A$, with its relative regular
neighbourhood homeomorphic to $B$.

Begin with collars of the components of $A$.  In a vertex-disc and
edge-band neighbourhood of the graph, absorb every additional vertex
disc along a forest rooted at those collars.  This leaves only surface
bands.  Since the initial collar has Euler characteristic $a$ and each
band decreases Euler characteristic by one, there are $a-\chi(B)$
bands.  Their incidence graph connects the $a+c$ initial collars.
The interval bundle over each band is a three-dimensional one-handle.
The interval bundle over an initial collar is merely an outward collar
of the given interface and is absorbed into the core.  For a circular
component, its loop preserves orientation in the base collar, so
orientability of the total space forces trivial fibre reversal around
it; its vertical interface is therefore an annulus with a product
collar.  The interval case is a product directly.

If $A=\partial B$, remove the interior of a small disc $\Delta$ in $B$.
The punctured surface $B^\circ$ has a free boundary circle and
$\chi(B^\circ)=\chi(B)-1$.  The preceding argument gives
$1-\chi(B)$ one-handles, since now $a=0$.  Restoring
$p^{-1}(\Delta)=\Delta\times I$ attaches one two-handle along
$\partial\Delta\times I$.  A possibly twisted surface band still
lifts to a three-dimensional one-handle; the orientability assumption
concerns the total space.  This also proves the nonorientable-base case.
\end{proof}

For a saturated base of orientable genus $g$ with $b$ boundary circles,
$h_1=2g+b-1$.  For a nonorientable base with $k$ crosscaps and $b$
boundary circles, $h_1=k+b-1$.  These can be large numerical quantities.
The argument gives an existence statement for the relative bands; it
does not require us to write a list with one entry for every band.

\subsection{The geometric residual incidence theorem}
\label{geom:residual}

\begin{theorem}[Residual incidence from the source annuli]
\label{geom:thm-residual-incidence}
Assume the geometric input of Section~\ref{geom:input}, including
compatibility with the original core cocores.  Let $E\to B$ be saturated,
with attaching-annulus midpoint curves $\gamma_1,\ldots,\gamma_b$.
There is a relative handle presentation with $r=1-\chi(B)$ new
one-handles and one residual two-handle $R$ such that:
\begin{enumerate}
\item Its attaching curve $\alpha_R$ has original-core incidence
\begin{equation}
 |\alpha_R\cap\partial D_h|
 =\sum_{j=1}^{b}|\gamma_j\cap\partial D_h|.
 \label{geom:eq-residual-row}
\end{equation}
\item It meets the cocore of each of these $r$ new one-handles exactly
twice.
\item Original core attaching curves and residual curves belonging to
other bundle components miss all these $r$ new cocores.
\end{enumerate}
All the attaching curves occur simultaneously as disjoint embedded
curves on the boundary of the enlarged zero/one-handle union.
\end{theorem}

\begin{proof}
Use the punctured surface $B^\circ=B\setminus\operatorname{int}\Delta$
from Lemma~\ref{geom:lem-relative-count}.  Its $b$ old boundary circles
are the attaching set, and its only free boundary circle is
$C=\partial\Delta$.  Choose its relative presentation as the regular
neighbourhood of the old boundary collars and $r$ bands.  The inner
edge of each collar remains one of the old boundary circles.  Every
other boundary segment must therefore lie on $C$.  In particular,
\emph{both long sides of every band occur on $C$, each exactly once}.

We arrange the entire band feet to avoid the old cocores.  First extend
the old cocore discs through the added interface collars by product
strips.  By Definition~\ref{geom:def-compatible}, their traces on each
vertical annulus are finitely many complete interval fibres.  On the
free edge of the corresponding base collar these fibres give a finite
set of marked points.

Shrink and relocate the finitely many band-foot intervals into small
unmarked arcs, preserving their cyclic order on that free edge.  An
orientation-preserving circle homeomorphism realizes this relocation.
Interpolate it across the collar to the identity on the \emph{actual
attaching edge} of $B$.  Thus the resulting collar isotopy fixes $A$
pointwise.  Its lift fixes the given attaching annulus pointwise and
does not modify the external gluing map.  Since the cocore traces are
entire fibres, the lifted foot rectangles and sufficiently small
neighbourhoods of them miss all old cocore boundaries.  Avoiding just
the midpoint points without this product hypothesis would not suffice.

The free circle $C$ follows the outer edge of each old collar except
where a band foot has replaced a short subarc.  Every old cocore
crossing lies outside those feet and is encountered exactly once.
The portions on the bands lie in new one-handles and add no traversal
of an original core one-handle.  This proves
\eqref{geom:eq-residual-row}.

Lift the bands to one-handles and restore $\Delta\times I$.  The
midpoint section of its attaching annulus gives the embedded curve
$\alpha_R$.  Each long side of a band crosses the lifted band cocore
once, so the two sides give exactly two geometric intersections.
Twisting a surface band can change the signs of these passages but
cannot change their unsigned number.

The original attaching curves lie outside the interface collars, and
different bundles have disjoint charts.  Their residual curves can
therefore be constructed disjointly, and no such curve uses another
bundle's bands.  The entire construction adds one-handles before
two-handles and creates no three-handle.  It is consequently a genuine
geometric handle presentation with all the asserted intersections.
\end{proof}

For $B$ a disc, $r=0$ and
\eqref{geom:eq-residual-row} is exactly the earlier disk-absorption
formula~\cite{marked-parallelity-85}.  For a higher-genus or nonorientable
base, the theorem compiles the same finite kind of row even though
the residual curve may run through a very large family of new handles.
It neither produces a short freely reduced word nor needs one.

\subsection{Connecting different cores and counting the private handles}
\label{geom:connectors}

Contract a spanning forest of the original zero/one-handle graph.
Let its connected components be indexed by $j$.  Component $j$ has
$v_j$ original zero-handles, $e_j$ original one-handles, and $f_j$
original two-handles.  Its remaining original one-handle count is
\[
 d_j=e_j-v_j+1.
\]
Contracting a tree handle deletes its incidence column and retains the
other original cocores.  It does not freely reduce any attaching curve.

Form the bipartite attachment multigraph $\Gamma$: its vertices are
the original core components and the attached bases, and each connected
interval or circle interface gives one edge.  Repeated attachments are
retained as separate edges.  For a connected component $C$ of this
graph, let $q_C$ be its number of core vertices and put
\begin{equation}
 H_C=\sum_{B_i\in C}(a_i-\chi(B_i)+\varepsilon_i),
 \qquad
 p_C=H_C-q_C+1.
 \label{geom:eq-private-count}
\end{equation}

\begin{lemma}[Connector forest]
\label{geom:lem-connector-forest}
The enlarged zero/one-handle union indexed by $C$ is connected.
Exactly $q_C-1$ of its added one-handles can be contracted as a
spanning tree joining the original core components.  Its remaining
added one-handles number $p_C\geq0$.
\end{lemma}

\begin{proof}
The relative-band graph for one attached base connects all its
interface collars.  After those collars are absorbed into the core,
it therefore connects precisely the core components incident to that
base.  Doing this for all bases gives the same connected-component
relation as $\Gamma$.  Higher-index handles do not join components.
Hence the added one-handle graph on the $q_C$ original core pieces is
connected.  Choose a spanning tree, which has $q_C-1$ edges, and
contract its corresponding handles.  The remaining number is
$H_C-(q_C-1)$, as claimed.
\end{proof}

We call these remaining added handles \emph{private} to distinguish
them from the original core columns on which the cancellation search
will operate.  Initially each private column has either only zero
entries or a single entry equal to two in its own saturated residual
row.  Tree contraction merely deletes some added columns and leaves
this property intact.

The correction by $q_C-1$ is essential.  A single saturated annulus
has $1-\chi(B)=1$ relative band.  If its two boundary annuli attach
to distinct core components, that band joins them and is contracted;
there is no surviving private handle.  More generally, one saturated
base meeting $s$ distinct cores leaves $1-\chi(B)-(s-1)$ added
handles after those mergers.  This is $2g+b-s$ for an orientable
base and $k+b-s$ for a nonorientable base.

\subsection{The finite active matrix and the handlebody census}
\label{geom:active-matrix}

The full geometric positive-pivot operation at an entry $n_{rg}=1$
has unsigned update rule
\begin{equation}
 n'_{ih}=n_{ih}+n_{ig}n_{rh}
 \qquad(i\ne r,\ h\ne g).
 \label{geom:eq-positive-pivot}
\end{equation}
The maintained geometric disc-slide theorem realizes this rule by
embedded disc slides followed by a one/two-handle cancellation
\cite{marked-parallelity-85}.  Its hypothesis is the actual one-point
intersection, not a unit algebraic intersection number.

An initially even column stays even under
\eqref{geom:eq-positive-pivot}.  More generally, an initial column
with entries only zero or at least two can never acquire a unit
entry; see Lemma~\ref{lem:inert-columns}.  Thus no private column
can ever be used as a pivot.  The original-column submatrix evolves
autonomously under the same rule, since an update on an original
column refers only to entries in that column and in the selected
original pivot column.

\begin{proposition}[Active incidence compilation]
\label{geom:prop-active-compilation}
After the forest contractions above, form a matrix $N$ as follows.
Keep the rows of the original core two-handles, restricted to the
surviving original one-handles.  Append one row for every saturated
base, with entries given by \eqref{geom:eq-residual-row}.  Do not
append any row for an unsaturated base and do not represent any
private column.

The full relative geometric presentation admits a positive-pivot
schedule eliminating every two-handle if and only if $N$ does.
Any successful schedule for $N$ certifies actual geometric
cancellations in the full presentation.  This equivalence remains
true when all positive entries are capped at two.
\end{proposition}

\begin{proof}
Private columns cannot pivot, and deleting them leaves the updates
on the other columns unchanged.  A full schedule therefore restricts
to a schedule on $N$.  Conversely, execute a schedule for $N$ on the
full matrix.  Its selected entries are still one, so the geometric
disc-slide and cancellation theorem applies at every step.  All rows
are removed.  The map $n\mapsto\min(2,n)$ preserves addition and
multiplication in the saturated nonnegative semiring, and detects
the predicate $n=1$ exactly.  It therefore preserves the argument.
\end{proof}

Theorem~\ref{thm:column-peeling} supplies the complete finite procedure
for full-row elimination in this raw model.  In particular, the
large matrix implicit in the relative presentation need never be
materialized: the producer and checker work with the original
core columns, the original rows, and one additional finite row per
saturated base.

\begin{theorem}[Certified whole-cut handlebody census]
\label{geom:thm-census}
Assume the source geometric contract, no core three-handles, and
bordered bases.  Suppose the compiled active matrix has a verified
schedule eliminating all its rows.  Then each component $X_C$
indexed by the attachment graph is a handlebody, of genus
\begin{align}
 G_C
 &=\sum_{j\in C}d_j+p_C-\sum_{j\in C}f_j
     -\sum_{i\in C}\varepsilon_i
 \label{geom:eq-census}\\
 &=\sum_{j\in C}(e_j-v_j+1-f_j)
   +\sum_{i\in C}(a_i-\chi(B_i))-q_C+1.
 \label{geom:eq-census-euler}
\end{align}
Every unported bordered bundle is a separate handlebody of genus
$1-\chi(B_i)$ and contributes its entire authenticated multiplicity.
These data determine the complete genus histogram of the cut.
\end{theorem}

\begin{proof}
After contracting the two forests, the connected zero/one-handle
union for $C$ is a handlebody with
$\sum_jd_j+p_C$ one-handles.  There are
$\sum_jf_j+\sum_i\varepsilon_i$ two-handles.  By
Proposition~\ref{geom:prop-active-compilation}, the checked schedule
cancels all of them geometrically, leaving exactly the number of
one-handles in \eqref{geom:eq-census}.  Substitution of
\eqref{geom:eq-private-count} gives
\eqref{geom:eq-census-euler}.

An orientable total interval bundle over a connected bordered
surface is a handlebody: choose proper arcs cutting the base to a
disc and lift them to cocore discs.  There are $1-\chi(B_i)$ arcs,
including for a nonorientable base.  An unported component has no
intersection with the core or with another bundle component, so it
must be counted separately with its full multiplicity.
\end{proof}

There is an independent Euler check.  The intersection of the core
with a proper interval interface is a disc, and with a complete
circle interface is an annulus.  Hence
\[
 \chi(X_C)=\sum_{j\in C}(v_j-e_j+f_j)
           +\sum_{i\in C}\chi(B_i)-\sum_{i\in C}a_i,
\]
which equals $1-G_C$ by \eqref{geom:eq-census-euler}.  This equality
checks the count; it does not substitute for the cancellation proof.
In particular, the conditional expression $G_C=0$ is not a ball
certificate when the active matrix has a surviving row.

\subsection{Examples and boundaries of the conclusion}
\label{geom:examples}

\paragraph{A saturated annulus joining two solid tori.}
Let its boundary annuli attach to solid-torus cores whose standard
meridians meet their midpoint curves $p$ and $q$ times respectively.
Choose the usual embedded slopes with these positive geometric
winding numbers.  The connecting band merges the two cores and
is contracted.  The active matrix is the single row $(p,q)$, with
no private columns.  If $p=1$ or $q=1$, one geometric cancellation
proves that the union is a solid torus.  The genus computation is
$2+0-1=1$.

For $p=2$ and $q=3$ there is no unit pivot.  This refusal is
necessary for this example: van Kampen gives
\[
 \langle x,y\mid x^2=y^3\rangle,
\]
which maps onto the nonabelian group $S_3$ by sending $x$ to a
transposition and $y$ to a three-cycle.  Its abelianization is
$\mathbb Z$, so first homology alone would not distinguish it
from the positive solid-torus cases.

\paragraph{A saturated M\"obius bundle.}
The orientable interval bundle over a M\"obius band is a solid
torus, with the vertical boundary representing the square of its
core generator.  Attach that annulus to a solid-torus core along
an annulus with one geometric meridian intersection.  The relative
description has one private one-handle and one residual row whose
active part is $(1)$.  Cancelling against the original core handle
leaves one private handle, proving a solid torus.  The corresponding
relation can be written $x=a^2$, but the positive conclusion comes
from the embedded cancellation.

If instead the same vertical annulus is attached to a
nullhomotopic annulus on the boundary of a ball, the active matrix
has one \emph{empty} row.  Its private column has incidence two.
No cancellation is authorized.  Indeed the resulting group is
$\langle a\mid a^2\rangle\cong\mathbb Z/2$, whereas its Euler
characteristic is one and its boundary is a sphere.  This is a
concrete reason to keep the surviving-row obligation separate
from the conditional genus arithmetic.

\paragraph{Arbitrarily large binary genus.}
Let $B$ be orientable of genus $g$ with one boundary component,
and attach its vertical boundary annulus to a solid-torus core
along a curve meeting a standard meridian once.  There are $2g$
private handles and one active unit entry.  The residual cancels
the original core handle, leaving a handlebody of genus $2g$.
For $g=2^L$ the genus counter requires $O(L)$ bits, while the
active matrix still contains just one entry.  This is a family
of explicitly specified relative topological attachments.  It
does not assert that every member has a normal-cut source
triangulation with $O(L)$ tetrahedra.

\paragraph{What successful augmentation does and does not imply.}
No core handlebody conclusion is assumed by the compiler, but a
successful full column-peeling schedule restricts to one for any
subset of its rows: deleting the other rows can only make its
pivot columns less occupied.  Thus successful joint cancellation
implies raw cancellability of the original core rows themselves.
Adding saturated residual rows cannot rescue a core outside that
raw cancellation class.  This observation also justifies an early
filter: an obstruction in the original core matrix permits this
consumer to stop before requesting the ordered-boundary queries.
It is a refusal by this consumer, not a proof of non-handlebodyness.

If the active matrix stalls, disc slides that reduce geometric
intersection numbers, another meridian system, or a different
handle presentation might still succeed.  Closed bases and core
three-handles remain outside the present operator.  A successful
histogram concerns the \emph{entire supplied cut}, including
isolated components.  A knot verdict additionally requires the
appropriate cutting-surface and peripheral information together
with authenticated provenance from the input knot diagram.

Finally, the private bands and the residual curve are constructed
by existence in the proof, while their required incidences are
compiled from source charts.  The certificate does not provide
a source address for an arbitrary private cocore, nor the
transport of a later boundary pattern through the chosen band
decomposition.  Closing those geometric interfaces, controlling
surface discovery, and bounding all hierarchy stages remain
necessary to reach the quasi-polynomial objective discussed in
Lackenby's program~\cite{lackenby2026}.
